\scanpage{161}
UNIVERSITE\\
DES SCIENCES ET TECHNIQUES\\
DU LANGUEDOC\\
------------

MONTPELLIER, le 16 septembre 1980

Service de la Scolarité\\
MA/IE\\
\uline{Diffusion} : liste A\\
N° 4583

\uline{AVIS DE SOUTENANCE D'UNE THESE DE DOCTORAT DE 3ème CYCLE} :

\uline{AGRONOMIE Option ZOOTECHNIE}

\uline{Nom et prénom du candidat} : JIMENEZ Nelys épouse PEREZ

\uline{Titre de la thèse} : Effets de la gestation et du mode de conduite (traite ou allaitement) sur la durée de l'anoestrus post-partum chez la vache française Frisonne Pie noire (FFPN).

\uline{Date, heure et lieu de soutenance prévus} : 26 septembre 1980 à 9 H 30, laboratoire de Zootechnie, ENSAM, Montpellier

\uline{Nom du Directeur de la thèse} : MM. DAUZIER et MAULEON

\uline{Désignation du laboratoire où la thèse est préparée} : laboratoire de Physiologie de la Reproduction - INRA - Nouzilly

\uline{Analyse de l'ouvrage} : Sur 80 vaches Frisonnes Pie noire (FFPN) une caractérisation hormonale pré et post-partum a été faite en comparant deux modes de conduite : traite ou allaitement.

L'importance du rôle de l'utérus dans la reprise de l'activité ovarienne a été \deleted{\unsure{le}} mise en évidence grâce à la mesure dans le sang du métabolite de la $\mathrm{PGF}_2\alpha$ (PGFM).

Le rôle déterminant du système nerveux dans cette reprise a été confirmé par l'étude objective des variations plasmatiques de LH.

Variations de concentration plasmatique de LH, niveau de prostaglandine F2$\alpha$ et prolactinémie expliquent 40 \% de la variabilité de l'intervalle vêlage-1ère ovulation. Des corrélations entre paramètres pre-partum et post-partum suggèrent qu'une partie des variations est issue des situations physiologiques pendant la gestation.

Les résultats montrent clairement l'effet \unsure{négatif|néggatif} sur la reprise de l'activité ovarienne de l'allaitement en comparaison avec la traite.

\uline{Durée approximative des recherches ayant abouti à la thèse} : 3 ans.

\uline{Composition du jury proposé} : MM. : DAUZIER\\
LEGRAND\\
MAULEON\\
VEZINHET

P/le Secrétaire Général\\
L'Attachée d'Administration Universitaire,

\figure{\begin{tikzpicture}[x=1cm,y=1cm]
\draw (0,0) .. controls (-.08,.25) and (.24,.5) .. (.19,.71) .. controls (.12,.83) and (.02,.31) .. (.07,.17) .. controls (.18,.01) and (.15,.37) .. (.01,.36) .. controls (-.08,.3) and (.03,.01) .. (.04,-.04);
\end{tikzpicture}}
\uline{\unsure{Gentil}}\\
\uline{GENTIL}

MR LE PROFESSEUR GROTHENDIECK ALEX.\\
MATHEMATIQUES\qquad5

\margin{161}
